\documentclass[12pt,a4paper]{report}

% ==================== PACKAGES ====================
\usepackage[utf8]{inputenc}
\usepackage[T1]{fontenc}
\usepackage[french]{babel}
\usepackage{graphicx}
\usepackage{geometry}
\usepackage{setspace}
\usepackage{hyperref}
\usepackage{titlesec}
\usepackage{float}
\usepackage{booktabs}
\usepackage{array}
\usepackage{xcolor}
\usepackage{enumitem}
\usepackage{fancyhdr}
\usepackage{longtable}
\usepackage{colortbl}
\usepackage{mdframed}
\usepackage{multicol}
\usepackage{verbatim}
\usepackage{multirow}

\geometry{margin=2.5cm}
\onehalfspacing

% ==================== COULEURS ====================
\definecolor{bleu3S}{RGB}{37, 99, 235}
\definecolor{rouge}{RGB}{220, 38, 38}
\definecolor{orangecol}{RGB}{249, 115, 22}
\definecolor{vertcol}{RGB}{34, 197, 94}
\definecolor{grisleger}{RGB}{243, 244, 246}

% ==================== EN-TETE/PIED DE PAGE ====================
\pagestyle{fancy}
\fancyhf{}
\rhead{\textcolor{bleu3S}{\textbf{Sprint 1}}}
\lhead{Backlog \& User Stories}
\cfoot{\thepage}
\renewcommand{\headrulewidth}{0.4pt}

% ==================== TITRES ====================
\titleformat{\chapter}[block]
  {\bfseries\LARGE\color{bleu3S}}
  {}{0em}{}
  [\vspace{0.2em}\titlerule]

\titleformat{\section}{\bfseries\large\color{bleu3S}}{\thesection}{1em}{}

% ==================== HYPERREF ====================
\hypersetup{colorlinks=true, linkcolor=bleu3S, urlcolor=bleu3S}

% ==================== DEBUT DOCUMENT ====================
\begin{document}

% ==================== PAGE DE GARDE ====================
\begin{titlepage}
    \begin{center}
        \vspace*{2cm}
        \textcolor{bleu3S}{\Large \textbf{Methodologie Agile Scrum}}\\[0.5cm]
        \rule{\textwidth}{0.4mm}\\[0.4cm]
        {\LARGE \textbf{Backlog Initial et User Stories}}\\[0.2cm]
        {\large \textbf{Sprint 1 -- Semaines 1 a 4}}\\[0.4cm]
        \rule{\textwidth}{0.4mm}
        \vspace{1.5cm}
        {\large \textbf{Projet :}}\\[0.2cm]
        {\Large 3S TalentMatch}\\
        {\large Plateforme Intelligente de Parsing de CVs}
        \vspace{2cm}
        \begin{minipage}{0.45\textwidth}
            \centering
            \textbf{Realise par :}\\[0.2cm]
            Youssef Gara\\
            Etudiant Ingenieur\\
            ESPRIT
        \end{minipage}
        \hfill
        \begin{minipage}{0.45\textwidth}
            \centering
            \textbf{Encadre par :}\\[0.2cm]
            Mme Sirine Nasri\\
            Mme Loujein\\[0.3cm]
            3S Group
        \end{minipage}
        \vspace{2cm}
        \textbf{Duree Sprint 1 :} 4 semaines (28 jours)\\
        \textbf{Date :} Fevrier 2026
    \end{center}
\end{titlepage}

\tableofcontents
\newpage

% ==================== INTRODUCTION ====================
\chapter*{Introduction}
\addcontentsline{toc}{chapter}{Introduction}

Ce document presente le backlog initial et les user stories detaillees du \textbf{Sprint 1} du projet 3S TalentMatch, conformement aux directives de l'encadrante.

\subsection*{Objectifs du Sprint 1}
\begin{enumerate}
    \item Finaliser et valider le cahier des charges
    \item Concevoir les diagrammes UML (cas d'utilisation, classes, sequence)
    \item Mettre en place l'environnement de developpement
    \item Implementer le MVP : Upload CV + Extraction texte (PDF/DOCX)
\end{enumerate}

\subsection*{Livrables Attendus}
\begin{itemize}
    \item Cahier des charges valide
    \item 3 diagrammes UML complets
    \item Backend FastAPI + PostgreSQL + Frontend React fonctionnels
    \item Demo operationnelle : Upload CV $\rightarrow$ Extraction $\rightarrow$ Affichage
\end{itemize}

% ==================== CHAPITRE 1 ====================
\chapter{Backlog Initial Priorise}

\section{Vue d'Ensemble}

Le backlog du Sprint 1 comprend \textbf{32 user stories} reparties en \textbf{6 epics}, pour un total de \textbf{113 story points}.

\begin{table}[H]
\centering
\begin{tabular}{|p{5cm}|p{2.5cm}|p{2cm}|p{3cm}|}
\hline
\rowcolor{bleu3S!20}
\textbf{Epic} & \textbf{Story Points} & \textbf{Priorite} & \textbf{Semaine} \\
\hline
Conception et Documentation & 19 & \textcolor{rouge}{\textbf{P0}} & Semaine 1 \\
\hline
Setup Environnement & 17 & \textcolor{rouge}{\textbf{P0}} & Semaines 1--2 \\
\hline
Backend Extraction CVs & 26 & \textcolor{rouge}{\textbf{P0}} & Semaines 2--3 \\
\hline
Base de Donnees & 19 & \textcolor{orangecol}{\textbf{P1}} & Semaines 2--3 \\
\hline
Frontend Upload Interface & 20 & \textcolor{orangecol}{\textbf{P1}} & Semaines 3--4 \\
\hline
Documentation et Tests & 12 & \textcolor{vertcol}{\textbf{P2}} & Semaine 4 \\
\hline
\rowcolor{grisleger}
\textbf{TOTAL} & \textbf{113} & -- & \textbf{4 semaines} \\
\hline
\end{tabular}
\caption{Resume des epics du Sprint 1}
\end{table}

\textbf{Legende :} \textcolor{rouge}{\textbf{P0}} Critique \quad \textcolor{orangecol}{\textbf{P1}} Important \quad \textcolor{vertcol}{\textbf{P2}} Souhaitable

\section{Epic 1 : Conception et Documentation}

\begin{longtable}{|p{1.5cm}|p{7cm}|p{2cm}|p{2cm}|}
\hline
\rowcolor{bleu3S!20}
\textbf{ID} & \textbf{User Story} & \textbf{Points} & \textbf{Priorite} \\
\hline
\endfirsthead
\hline \rowcolor{bleu3S!20} \textbf{ID} & \textbf{User Story} & \textbf{Points} & \textbf{Priorite} \\ \hline
\endhead
US-001 & Finalisation cahier des charges & 5 & P0 \\ \hline
US-002 & Diagramme de cas d'utilisation & 3 & P0 \\ \hline
US-003 & Diagramme de classes & 5 & P0 \\ \hline
US-004 & Diagramme de sequence (Upload CV) & 3 & P0 \\ \hline
US-005 & Documentation architecture technique & 3 & P1 \\ \hline
\rowcolor{grisleger} \textbf{Total Epic 1} & & \textbf{19} & \\ \hline
\caption{User Stories -- Epic 1}
\end{longtable}

\section{Epic 2 : Setup Environnement}

\begin{longtable}{|p{1.5cm}|p{7cm}|p{2cm}|p{2cm}|}
\hline
\rowcolor{bleu3S!20}
\textbf{ID} & \textbf{User Story} & \textbf{Points} & \textbf{Priorite} \\
\hline
\endfirsthead
\hline \rowcolor{bleu3S!20} \textbf{ID} & \textbf{User Story} & \textbf{Points} & \textbf{Priorite} \\ \hline
\endhead
US-006 & Configuration environnement Python (venv) & 2 & P0 \\ \hline
US-007 & Installation FastAPI + dependances & 2 & P0 \\ \hline
US-008 & Configuration PostgreSQL local & 3 & P0 \\ \hline
US-009 & Setup projet React avec Vite & 3 & P1 \\ \hline
US-010 & Configuration Git + GitHub repository & 2 & P1 \\ \hline
US-011 & Docker Compose (backend + BDD) & 5 & P2 \\ \hline
\rowcolor{grisleger} \textbf{Total Epic 2} & & \textbf{17} & \\ \hline
\caption{User Stories -- Epic 2}
\end{longtable}

\section{Epic 3 : Backend -- Upload et Extraction CVs}

\begin{longtable}{|p{1.5cm}|p{7cm}|p{2cm}|p{2cm}|}
\hline
\rowcolor{bleu3S!20}
\textbf{ID} & \textbf{User Story} & \textbf{Points} & \textbf{Priorite} \\
\hline
\endfirsthead
\hline \rowcolor{bleu3S!20} \textbf{ID} & \textbf{User Story} & \textbf{Points} & \textbf{Priorite} \\ \hline
\endhead
US-012 & API endpoint POST /upload-cv & 5 & P0 \\ \hline
US-013 & Extraction texte PDF avec PyPDF2 & 5 & P0 \\ \hline
US-014 & Extraction texte DOCX avec python-docx & 3 & P0 \\ \hline
US-015 & Detection automatique type fichier & 3 & P0 \\ \hline
US-016 & Validation format fichier (taille, type) & 2 & P1 \\ \hline
US-017 & Gestion erreurs extraction & 3 & P1 \\ \hline
US-018 & Tests unitaires extraction (Pytest) & 5 & P1 \\ \hline
\rowcolor{grisleger} \textbf{Total Epic 3} & & \textbf{26} & \\ \hline
\caption{User Stories -- Epic 3}
\end{longtable}

\section{Epic 4 : Base de Donnees}

\begin{longtable}{|p{1.5cm}|p{7cm}|p{2cm}|p{2cm}|}
\hline
\rowcolor{bleu3S!20}
\textbf{ID} & \textbf{User Story} & \textbf{Points} & \textbf{Priorite} \\
\hline
\endfirsthead
\hline \rowcolor{bleu3S!20} \textbf{ID} & \textbf{User Story} & \textbf{Points} & \textbf{Priorite} \\ \hline
\endhead
US-019 & Creation schema BDD (tables initiales) & 5 & P0 \\ \hline
US-020 & Table users (authentification) & 3 & P1 \\ \hline
US-021 & Table candidates (infos candidats) & 3 & P1 \\ \hline
US-022 & Table cv\_documents (metadonnees CVs) & 3 & P1 \\ \hline
US-023 & SQLAlchemy models + migrations Alembic & 5 & P1 \\ \hline
\rowcolor{grisleger} \textbf{Total Epic 4} & & \textbf{19} & \\ \hline
\caption{User Stories -- Epic 4}
\end{longtable}

\section{Epic 5 : Frontend -- Interface Upload}

\begin{longtable}{|p{1.5cm}|p{7cm}|p{2cm}|p{2cm}|}
\hline
\rowcolor{bleu3S!20}
\textbf{ID} & \textbf{User Story} & \textbf{Points} & \textbf{Priorite} \\
\hline
\endfirsthead
\hline \rowcolor{bleu3S!20} \textbf{ID} & \textbf{User Story} & \textbf{Points} & \textbf{Priorite} \\ \hline
\endhead
US-024 & Page upload avec drag \& drop & 8 & P0 \\ \hline
US-025 & Affichage texte extrait & 3 & P0 \\ \hline
US-026 & Gestion etats chargement (loading) & 2 & P1 \\ \hline
US-027 & Affichage erreurs utilisateur & 2 & P1 \\ \hline
US-028 & Design responsive (mobile/desktop) & 5 & P2 \\ \hline
\rowcolor{grisleger} \textbf{Total Epic 5} & & \textbf{20} & \\ \hline
\caption{User Stories -- Epic 5}
\end{longtable}

\section{Epic 6 : Documentation et Tests}

\begin{longtable}{|p{1.5cm}|p{7cm}|p{2cm}|p{2cm}|}
\hline
\rowcolor{bleu3S!20}
\textbf{ID} & \textbf{User Story} & \textbf{Points} & \textbf{Priorite} \\
\hline
\endfirsthead
\hline \rowcolor{bleu3S!20} \textbf{ID} & \textbf{User Story} & \textbf{Points} & \textbf{Priorite} \\ \hline
\endhead
US-029 & README backend avec instructions & 3 & P1 \\ \hline
US-030 & README frontend avec instructions & 2 & P1 \\ \hline
US-031 & Documentation API (Swagger) & 2 & P2 \\ \hline
US-032 & Tests E2E (upload vers extraction) & 5 & P2 \\ \hline
\rowcolor{grisleger} \textbf{Total Epic 6} & & \textbf{12} & \\ \hline
\caption{User Stories -- Epic 6}
\end{longtable}

% ==================== CHAPITRE 2 ====================
\chapter{User Stories Detaillees}

\section{Format Standard}

\begin{mdframed}[backgroundcolor=grisleger!50,linecolor=bleu3S,linewidth=1pt]
\textbf{En tant que} [acteur] \\
\textbf{Je veux} [fonctionnalite] \\
\textbf{Afin de} [benefice/objectif]
\end{mdframed}

\section{US-001 : Finalisation Cahier des Charges}

\textbf{En tant que} etudiant ---
\textbf{Je veux} finaliser et faire valider le CDC ---
\textbf{Afin de} avoir une reference documentaire complete

\subsection*{Criteres d'Acceptation}
\begin{itemize}
    \item[\checkmark] Toutes sections obligatoires presentes (contexte, fonctionnel, architecture, IA, RGPD, UML, planning)
    \item[\checkmark] Section IA : approche CPU, inférence CamemBERT, pas de GPU ni fine-tuning
    \item[\checkmark] Document compile en PDF sans erreur LaTeX
    \item[\checkmark] Valide par Mme Nasri avant la fin de la Semaine 1
\end{itemize}

\textbf{Story Points : 5} \hfill \textbf{Priorite : P0}

\section{US-002 : Diagramme de Cas d'Utilisation}

\textbf{En tant que} chef de projet ---
\textbf{Je veux} creer un diagramme de cas d'utilisation ---
\textbf{Afin de} identifier tous les acteurs et leurs interactions

\subsection*{Criteres d'Acceptation}
\begin{itemize}
    \item[\checkmark] 3 acteurs identifies : Candidat, Recruteur RH, Administrateur
    \item[\checkmark] Minimum 13 cas d'utilisation (voir liste ci-dessous)
    \item[\checkmark] Diagramme cree avec Draw.io ou StarUML
    \item[\checkmark] Exporte en PNG et integre au CDC
\end{itemize}

\subsection*{Cas d'Utilisation}
\begin{multicols}{2}
\begin{enumerate}
    \item Creer compte
    \item Se connecter
    \item Uploader CV
    \item Consulter profil extrait
    \item Modifier donnees
    \item Creer offre emploi
    \item Consulter candidatures
    \item Lancer matching CV-Offre
    \item Consulter dashboard
    \item Exporter rapport
    \item Gerer utilisateurs
    \item Consulter logs RGPD
    \item Supprimer donnees
\end{enumerate}
\end{multicols}

\textbf{Story Points : 3} \hfill \textbf{Priorite : P0}

\section{US-003 : Diagramme de Classes}

\textbf{En tant que} developpeur backend ---
\textbf{Je veux} concevoir un diagramme de classes ---
\textbf{Afin de} definir le modele de donnees

\subsection*{Classes Principales}

\begin{table}[H]
\centering
\small
\begin{tabular}{|p{3cm}|p{10cm}|}
\hline
\rowcolor{bleu3S!20}
\textbf{Classe} & \textbf{Attributs et Methodes} \\
\hline
User & \texttt{user\_id, email, password\_hash, role} / \texttt{authenticate()} \\
\hline
Candidat & \texttt{candidat\_id, nom, email, telephone} / \texttt{upload\_cv()} \\
\hline
CVDocument & \texttt{cv\_id, file\_path, extracted\_text, created\_at} \\
\hline
CVData & \texttt{structured\_data, competences, experiences} / \texttt{to\_json()} \\
\hline
Offre & \texttt{offre\_id, titre, competences\_requises} / \texttt{match\_candidats()} \\
\hline
Matching & \texttt{matching\_id, score} / \texttt{calculate()} \\
\hline
AccessLog & \texttt{log\_id, user\_id, action, timestamp} / \texttt{log()} \\
\hline
Consentement & \texttt{consent\_id, stockage, analyse\_ia} / \texttt{revoke()} \\
\hline
\end{tabular}
\caption{Classes principales du systeme}
\end{table}

\textbf{Story Points : 5} \hfill \textbf{Priorite : P0}

\section{US-012 : API Endpoint POST /upload-cv}

\textbf{En tant que} developpeur backend ---
\textbf{Je veux} creer un endpoint d'upload de CV ---
\textbf{Afin de} recevoir des fichiers depuis le frontend

\subsection*{Criteres d'Acceptation}
\begin{itemize}
    \item[\checkmark] \texttt{POST /api/upload-cv} accepte multipart/form-data
    \item[\checkmark] Validation : taille max 10 MB, formats PDF/DOCX uniquement
    \item[\checkmark] Retourne JSON avec \texttt{cv\_id}, \texttt{message}, \texttt{success}
    \item[\checkmark] Gestion erreurs HTTP 400, 413, 500
    \item[\checkmark] Documentation Swagger generee sur \texttt{/docs}
\end{itemize}

\subsection*{Exemple de Code}
\begin{verbatim}
from fastapi import APIRouter, UploadFile, File, HTTPException
import uuid

router = APIRouter(prefix="/api")

@router.post("/upload-cv")
async def upload_cv(file: UploadFile = File(...)):
    contents = await file.read()
    if len(contents) > 10 * 1024 * 1024:
        raise HTTPException(413, "Fichier trop volumineux")
    cv_id = str(uuid.uuid4())
    return {"success": True, "cv_id": cv_id}
\end{verbatim}

\textbf{Story Points : 5} \hfill \textbf{Priorite : P0}

\section{US-013 : Extraction Texte PDF}

\textbf{En tant que} systeme ---
\textbf{Je veux} extraire le texte d'un PDF textuel ---
\textbf{Afin de} l'envoyer au pipeline NLP

\subsection*{Exemple de Code}
\begin{verbatim}
from pypdf import PdfReader

def extract_text_from_pdf(file_path: str) -> dict:
    reader = PdfReader(file_path)
    text = "\n\n".join(
        p.extract_text() for p in reader.pages
        if p.extract_text()
    )
    if len(text.strip()) < 100:
        return {"success": False,
                "error": "PDF scanne, OCR necessaire"}
    return {"success": True, "text": text}
\end{verbatim}

\textbf{Story Points : 5} \hfill \textbf{Priorite : P0}

\section{US-024 : Page Upload avec Drag \& Drop}

\textbf{En tant que} recruteur ---
\textbf{Je veux} uploader un CV par drag \& drop ---
\textbf{Afin de} soumettre facilement un CV

\subsection*{Criteres d'Acceptation}
\begin{itemize}
    \item[\checkmark] Zone drag \& drop avec feedback visuel au survol
    \item[\checkmark] Bouton ``Parcourir'' alternatif
    \item[\checkmark] Loading spinner pendant upload
    \item[\checkmark] Message succes ou erreur apres upload
    \item[\checkmark] Affichage du texte extrait dans l'interface
\end{itemize}

\textbf{Story Points : 8} \hfill \textbf{Priorite : P0}

% ==================== CHAPITRE 3 ====================
\chapter{Planning Sprint 1}

\begin{table}[H]
\centering
\small
\begin{tabular}{|p{1.5cm}|p{2cm}|p{2cm}|p{4cm}|p{4cm}|}
\hline
\rowcolor{bleu3S!20}
\textbf{Sem.} & \textbf{Jour} & \textbf{US} & \textbf{Activites} & \textbf{Livrable} \\
\hline
\multirow{5}{*}{\textbf{S1}} & Lundi & US-001, 002 & CDC + Diag. cas utilisation & CDC PDF \\
\cline{2-5}
 & Mardi & US-003 & Diagramme de classes & Classes PNG \\
\cline{2-5}
 & Mercredi & US-004, 005 & Diag. sequence + Doc archi & Diagrammes \\
\cline{2-5}
 & Jeudi & US-006, 007 & Setup Python + FastAPI & Backend init \\
\cline{2-5}
 & Vendredi & US-008, 010 & PostgreSQL + Git & BDD + repo \\
\hline
\multirow{5}{*}{\textbf{S2}} & Lundi & US-012 & Endpoint /upload-cv & API upload \\
\cline{2-5}
 & Mardi & US-013 & Extraction PDF & pdf\_extractor.py \\
\cline{2-5}
 & Mercredi & US-014, 015 & Extraction DOCX + detection & docx\_extractor.py \\
\cline{2-5}
 & Jeudi & US-019, 021 & Schema BDD + Table candidates & Migration Alembic \\
\cline{2-5}
 & Vendredi & US-023 & SQLAlchemy models & models.py \\
\hline
\multirow{5}{*}{\textbf{S3}} & Lundi & US-016, 017 & Validations + erreurs & Handlers \\
\cline{2-5}
 & Mardi & US-018 & Tests unitaires Pytest & test\_*.py \\
\cline{2-5}
 & Mercredi & -- & Buffer / Debugging & -- \\
\cline{2-5}
 & Jeudi & US-009 & Setup React + Vite & Frontend init \\
\cline{2-5}
 & Vendredi & US-024 debut & Page upload structure & UploadPage.tsx \\
\hline
\multirow{5}{*}{\textbf{S4}} & Lundi & US-024 suite & Drag \& drop + upload & Component complet \\
\cline{2-5}
 & Mardi & US-025, 026 & Affichage resultat + loading & UI resultats \\
\cline{2-5}
 & Mercredi & US-027, 028 & Erreurs + responsive & UI finalisee \\
\cline{2-5}
 & Jeudi & US-029, 030 & READMEs & Documentation \\
\cline{2-5}
 & Vendredi & US-032 + Demo & Test E2E + Demo encadrantes & Demo Sprint 1 \\
\hline
\end{tabular}
\caption{Planning jour par jour du Sprint 1}
\end{table}

% ==================== CHAPITRE 4 ====================
\chapter{Definition de Termine et Metriques}

\section{Definition de ``Fini'' (DoD)}

\begin{itemize}
    \item[\checkmark] Code ecrit, fonctionne et commente
    \item[\checkmark] Tests unitaires passes (couverture > 70\% modules critiques)
    \item[\checkmark] Tests manuels valides
    \item[\checkmark] Code merge dans la branche \texttt{develop} sur GitHub
    \item[\checkmark] Criteres d'acceptation de la US valides
    \item[\checkmark] Demo possible aux encadrantes si applicable
\end{itemize}

\section{Velocite et Risques}

\textbf{Capacite Sprint 1 :} 80--100 points (etudiant solo). Total backlog : 113 points.\\
\textbf{Strategie :} Prioriser P0 et P1, reporter les US P2 au Sprint 2 si retard.

\begin{table}[H]
\centering
\begin{tabular}{|p{4cm}|p{2cm}|p{2cm}|p{5cm}|}
\hline
\rowcolor{bleu3S!20}
\textbf{Risque} & \textbf{Proba.} & \textbf{Impact} & \textbf{Mitigation} \\
\hline
Retard validation CDC & Moyenne & Eleve & Envoyer semaine 1 \\
\hline
Surcharge story points & Elevee & Moyen & Reporter P2 au Sprint 2 \\
\hline
Blocage technique & Faible & Moyen & Docs officielles + encadrantes \\
\hline
PostgreSQL local & Faible & Moyen & Docker Compose en backup \\
\hline
\end{tabular}
\caption{Matrice des risques}
\end{table}

% ==================== CONCLUSION ====================
\chapter*{Conclusion}
\addcontentsline{toc}{chapter}{Conclusion}

\begin{itemize}
    \item \textbf{32 user stories} detaillees avec criteres d'acceptation
    \item \textbf{113 story points} sur 4 semaines
    \item \textbf{6 epics} priorises (P0/P1/P2)
    \item \textbf{Planning jour par jour} sur 4 semaines
    \item \textbf{Definition de ``Fini''} claire et applicable
\end{itemize}

\subsection*{Prochaines Etapes}
\begin{enumerate}
    \item Validation de ce backlog par les encadrantes
    \item Creation du board Trello (voir guide)
    \item Demarrage Sprint 1 Semaine 1
    \item Daily standup auto-suivi chaque matin
    \item Revue de sprint fin de Semaine 4
\end{enumerate}

\vspace{2cm}
\begin{center}
Youssef Gara --- Fevrier 2026
\end{center}

\end{document}
